% !TeX root = ../../MainBook.tex
\section{分布与序列的谱密度}\label{sec:Discrete_Distribution}

如何分析$l^1$以外的序列，例如$x[n]=\cos(\Omega_0 n)$的频率成分？

当$\Omega_0/\pi$是无理数时，$x[n]$不是周期序列，不能对其做DFT。一种直接的想法是利用z变换。然而，\refsub{dtft}中指出，序列频谱在z变换的收敛域包含单位圆的情况下就是z变换在单位圆上的取值，当单位圆不在z变换的收敛域中时，这一方法就失效了——对于$l^1$以外的序列，z变换仅仅是一种变换域分析的手段，无法展示序列的频率成分。例如，对于$x[n]=\cos(\Omega_0 n)$，根据\cref{def:ROC}，$r=R=1$，$x[n]$的z变换没有收敛域。

可见，目前为止，我们还没有一种合适的工具来分析此类序列，包括最基本的复指数序列和正余弦序列。直观上，我们需要引入冲激项$\delta$来描述这些序列的频谱，需要做的工作与\refsec{Distributions}是类似的。

\subsection{离散时间傅里叶变换的分布理论}\label{sub:DTFT_distribution}

\begin{definition}[测试函数]
    在DTFT中，记圆周$\mathbb{T}=\mathbb{R}/2\pi\mathbb{Z}$（即周期为$2\pi$的实轴商空间），测试函数空间$\mathcal{C}(\mathbb{T})$为周期为$2\pi$的光滑函数（即无限阶可导的函数），其中的元素称为\textbf{测试函数}。
\end{definition}

\begin{definition}[速降序列]
    如果序列$\phi[n]$的衰减速度比任意多项式更快：
    \begin{equation}
        \lim_{|n|\to\infty}|n|^k\phi[n]=0,\quad \forall k\in\mathbb{N},
    \end{equation}
    则称$\phi[n]$是速降序列，将速降序列空间记为$\mathcal{S}(\mathbb{Z})$。
\end{definition}

\begin{proposition}
    设$\Phi(\Omega)$是周期为$2\pi$的函数，其傅里叶系数记为$\phi[n]$，则：
    \begin{equation}
        \phi[n]\in\mathcal{S}(\mathbb{Z})\iff \Phi(\Omega)\in\mathcal{C}(\mathbb{T}).
    \end{equation}
\end{proposition}
证明此命题需要用到附录中的结论。
\begin{proof}
    一方面，当$\phi[n]\in\mathcal{S}(\mathbb{Z})$时，DTFT满足
    \begin{equation}
        \abs{\Phi(\Omega)}=\abs{\sum_{n=-\infty}^{\infty}\phi[n]\e^{\im\Omega n}}\leq \sum_{n=-\infty}^{\infty}|\phi[n]|.
    \end{equation}
    容易由$\phi[n]\in\mathcal{S}(\mathbb{Z})$的定义得到$\sum_{n=-\infty}^{\infty}|\phi[n]|<\infty$，根据魏尔斯特拉斯M判别法，$\sum_{n=-\infty}^{\infty}\phi[n]\e^{\im\Omega n}$是一致收敛的函数项级数，逐项求导所得的函数项级数为：
    \begin{equation}
        \sum_{n=-\infty}^{\infty}\phi[n](\e^{\im\Omega n})'=\im\sum_{n=-\infty}^{\infty}n\phi[n]\e^{\im\Omega n}.
    \end{equation}
    类似可得：
    \begin{equation}
        \phi[n]\in\mathcal{S}(\mathbb{Z})\implies n\phi[n]\in\mathcal{S}(\mathbb{Z})\implies \im\sum_{n=-\infty}^{\infty}n\phi[n]\e^{\im\Omega n}\quad \text{一致收敛}.
    \end{equation}
    因此$\Phi(\Omega)$可导，并且
    \begin{equation}
        \Phi'(\Omega)=\im\sum_{n=-\infty}^{\infty}n\phi[n]\e^{\im\Omega n}.
    \end{equation}
    可以归纳地重复上述过程，最终得到$\Phi(\Omega)\in\mathcal{C}(\mathbb{T})$。

    另一方面，根据\cref{thm:converge_speed}，$\Phi(\Omega)\in\mathcal{C}(\mathbb{T})\implies\phi[n]\in\mathcal{S}(\mathbb{Z})$。
\end{proof}

\begin{definition}[缓增序列]
    如果存在常数$C$和$k\in\mathbb{N}$，使得
    \begin{equation}
        |x[n]|\leq C(1+|n|)^k,\quad \forall n\in\mathbb{Z},
    \end{equation}
    则称$x[n]$是\textbf{缓增序列}，将缓增序列空间记为$\mathcal{T}(\mathbb{Z})$。
\end{definition}

使用该定义描述缓增序列是合理的：$(1+|n|)^k$中，$1$放宽了对序列在$|n|$较小时的要求，而$|n|^k$描述序列递增速度不超过$k$次的多项式。

$\mathcal{C}(\mathbb{T})$的对偶空间是分布空间$\mathcal{D}(\mathbb{T})$，速降序列空间$\mathcal{S}(\mathbb{Z})$的对偶空间是缓增序列空间$\mathcal{T}(\mathbb{Z})$。

下面定义时域信号的配对和频域信号的配对。设$x[n],\phi[n]\in l^1\cap l^2$，$X(\e^{\im\Omega})$为$x[n]$的DTFT，$\phi[n]$为$\Phi(\Omega)$的傅里叶系数，则
\begin{equation}
    \IDTFT(\Phi(\Omega))[n]=\frac{1}{2\pi}\int_{-\pi}^{\pi}\Phi(\Omega)\e^{\im\Omega n}\,\dd\Omega=\phi[-n].
\end{equation}
在\cref{thm:DTFT_plancherel}中取$y[n]=\overline{\phi[n]}$，则$Y(\e^{\im\Omega})=\overline{\Phi(\Omega)}$，得到：
\begin{equation}
    \sum_{n=-\infty}^{\infty}x[n]\phi[n]=\frac{1}{2\pi}\int_{-\pi}^{\pi}X(\e^{\im\Omega})\Phi(\Omega)\,\dd\Omega.
\end{equation}

对于分布做同样的定义：

\begin{definition}[配对]
    时域信号的\textbf{配对}(pairing)定义为：
    \begin{equation}
        \langle x,\phi\rangle=\sum_{n=-\infty}^{\infty}x[n]\phi[n],\quad x\in\mathcal{T}(\mathbb{Z}),\ \phi\in\mathcal{S}(\mathbb{Z}),
    \end{equation}
    其中$x[n]$为缓增序列而$\phi[n]$为速降序列。

    数字角频域信号的配对定义为：
    \begin{equation}
        \langle X,\Phi\rangle=\frac{1}{2\pi}\int_{-\pi}^{\pi}X(\e^{\im\Omega})\Phi(\Omega)\,\dd\Omega,\quad X\in\mathcal{D}(\mathbb{T}),\ \Phi\in\mathcal{C}(\mathbb{T}).
    \end{equation}
\end{definition}

\begin{definition}[缓增序列的离散时间傅里叶变换]
    对于缓增序列$x[n]$，称分布$X$是$x[n]$的DTFT，如果：
    \begin{equation}
        \langle X,\Phi\rangle=\langle x,\phi\rangle,\quad \forall \phi\in\mathcal{S}(\mathbb{Z}),
    \end{equation}
    其中$\phi[n]$是$\Phi\in\mathcal{C}(\mathbb{T})$的傅里叶系数。
\end{definition}

\begin{example}[复指数序列的离散时间傅里叶变换]\label{exmp:exp_DTFT}
    求复指数序列$x[n]=\e^{\im\Omega_0 n}\ (-\pi<\Omega_0\leq\pi)$的DTFT。

    解：记$x[n]$的DTFT为$X$，由定义，
    \begin{equation}
        \langle X,\Phi\rangle=\langle x,\phi\rangle=\sum_{n=-\infty}^{\infty}\e^{\im\Omega_0 n}\phi[n],\quad \forall \phi\in\mathcal{S}(\mathbb{Z}).
    \end{equation}
    光滑函数的傅里叶级数一致收敛于自身（\cref{thm:uniform_convergency}），因此
    \begin{equation}
        \Phi(\Omega)=\sum_{n=-\infty}^{\infty}\phi[n]\e^{\im\Omega n}.
    \end{equation}
    可见：（注意数字角频域配对中的系数$2\pi$）
    \begin{equation}
        \langle X,\Phi\rangle=\sum_{n=-\infty}^{\infty}\e^{\im\Omega_0 n}\phi[n]=\Phi(\e^{\im\Omega_0})=\langle 2\pi\delta_{\Omega_0},\Phi\rangle.
    \end{equation}
    数字角频域总是以$2\pi$为周期的，根据以往的经验，$x[n]$的DTFT应该是以下周期分布：
    \begin{equation}
        X=2\pi\sum_{k=-\infty}^{\infty}\delta_{\Omega_0+2\pi k}=2\pi\tau_{\Omega_0}\shah_{2\pi},
    \end{equation}
    其中，平移算子$\tau$的定义可以参考\cref{def:distribution_delay}，$\shah_{2\pi}$是\cref{def:shah}给出的以$2\pi$为周期的狄拉克梳状分布。

    但是，事实上，我们不需要采用像$2\pi\tau_{\Omega_0}\shah_{2\pi}$这么复杂的符号，这一方面是因为$2\pi\delta_{\Omega_0}$已经足以描述$X$的行为，并且其意义是不言自明的；另一方面，从本节中对于配对的定义可以看出，$\mathcal{C}(\mathbb{T}),\mathcal{D}(\mathbb{T})$是否进行周期化并不影响分布的运算，换言之，$2\pi\delta_{\Omega_0}$和$2\pi\tau_{\Omega_0}\shah_{2\pi}$在$\mathcal{D}(\mathbb{T})$中\textbf{完全是同一个分布}。因此，$X=2\pi\delta_{\Omega_0}$。
\end{example}

作为$\Omega_0=0$时的特例，常数序列$x[n]=1$的DTFT为$2\pi\delta$。

在开始讨论分布理论中的DTFT的性质之前，我们仿照\refsec{Distributions}给出以下定义。

\begin{definition}[分布与函数的乘法]
    设$X\in\mathcal{D}(\mathbb{T}),\ g,\Phi\in\mathcal{C}(\mathbb{T})$，则$g$与$X$的乘积$gX$由下式定义：
    \begin{equation}
        \langle gX,\Phi\rangle=\langle X,g\Phi\rangle.
    \end{equation}
\end{definition}

\begin{definition}[分布的微分]
    设$X\in\mathcal{D}(\mathbb{T}),\ \Phi\in\mathcal{C}(\mathbb{T})$，则$X$的微分$X'$由下式定义：
    \begin{equation}
        \langle X',\Phi\rangle=-\langle X,\Phi'\rangle.
    \end{equation}
\end{definition}

\begin{definition}[分布的平移与伸缩]
    设$X\in\mathcal{D}(\mathbb{T}),\ \Phi\in\mathcal{C}(\mathbb{T})$，则$X$向右平移$\Omega_0$所得分布记为$\tau_{\Omega_0}X$，满足：
    \begin{equation}
        \langle \tau_{\Omega_0}X,\Phi\rangle=\langle X,\tau_{-\Omega_0}\Phi\rangle.
    \end{equation}
    由于配对是有限区间上的积分，不能定义$\sigma_a X$，但可以定义$X$的翻转，记作$X^-$：
    \begin{equation}
        \langle X^-,\Phi\rangle=\langle X,\Phi^-\rangle.
    \end{equation}
\end{definition}

\begin{definition}[分布的共轭]
    设$X\in\mathcal{D}(\mathbb{T}),\ \Phi\in\mathcal{C}(\mathbb{T})$，则$X$的共轭$\overline{X}$由下式定义：
    \begin{equation}
        \langle \overline{X},\Phi\rangle=\overline{\langle X,\overline{\Phi}\rangle}.
    \end{equation}
    由此可以定义分布的实部与虚部：
    \begin{equation}
        \Re(X)=\frac{X+\overline{X}}{2},\quad \Im(X)=\frac{X-\overline{X}}{2\im}.
    \end{equation}
\end{definition}

\begin{definition}[分布的卷积]
    设$X,Y\in\mathcal{D}(\mathbb{T}),\ \Phi\in\mathcal{C}(\mathbb{T})$，则$X,Y$的周期卷积为
    \begin{equation}
        \langle X*Y,\Phi\rangle=\langle X(u),\langle Y(v),\Phi(u+v)\rangle\rangle.
    \end{equation}
\end{definition}
与\cref{def:conv_distr_distr}不同，由于是有限区间上的积分，分布的卷积总是有定义的。%与\cref{exmp:associative_loss}的比较？

请读者自行验证上述定义对于正则分布成立。周期函数的卷积见\cref{def:period_conv}。

容易验证，对于时域的配对而言，同样有：
\begin{equation}
    \langle gx,\phi\rangle=\langle x,g\phi\rangle,\quad \langle \tau_{n_0}x,\phi\rangle=\langle x,\tau_{-n_0}\phi\rangle,\quad \langle \overline{x},\phi\rangle=\overline{\langle x,\overline{\phi}\rangle},\quad x,\phi,g\in\mathcal{S}(\mathbb{Z}),\ n_0\in\mathbb{Z}.
\end{equation}

此外，设$x\in\mathcal{T}(\mathbb{Z}),\ \phi\in\mathcal{S}(\mathbb{Z})$，考虑将时域卷积写成配对的形式：
\begin{equation}
    \langle x*y,\phi\rangle=\langle x,(y^-)*\phi\rangle,
\end{equation}
此时我们希望
\begin{equation}
    x*y\in\mathcal{T}(\mathbb{Z}),\quad (y^-)*\phi\in\mathcal{S}(\mathbb{Z}).
\end{equation}
与模拟信号的情况一样，可以有多种条件用来保证这一点，其中一种条件为$y\in\mathcal{S}(\mathbb{Z})$，这也是我们之后主要考虑的情况。

\subsection{离散时间傅里叶变换的分布的运算性质}\label{sub:DTFT_distribution_operation}

本小节中依次验证\refsub{property_of_dtft}中的DTFT的性质。不做特殊说明时默认：
\begin{equation}
    x[n],y[n]\in\mathcal{T}(\mathbb{Z}),\quad \phi[n]\in\mathcal{S}(\mathbb{Z}),
\end{equation}
并用符号$X,Y\in\mathcal{D}(\mathbb{T})$表示$x,y$的DTFT，$\phi[n]$是$\Phi(\Omega)$的傅里叶系数，即用$\Phi(\Omega)\in\mathcal{C}(\mathbb{T})$表示$\phi^-$的DTFT。

\begin{proposition}[线性性]
    \begin{equation}
        \DTFT(ax[n]+by[n])=aX+bY,\quad \forall a,b\in\mathbb{C}.
    \end{equation}
\end{proposition}
该命题直接得自配对$\langle\cdot,\cdot\rangle$对两个变元的线性性。

\begin{proposition}[时移性质]
    \begin{equation}
        \DTFT(x[n-n_0])=\e^{-\im \Omega n_0}X.
    \end{equation}
\end{proposition}
\begin{proof}
    \begin{equation}
        \phi[n+n_0]=\frac{1}{2\pi}\int_{-\pi}^{\pi}\Phi(\Omega)\e^{-\im\Omega (n+n_0)}\,\dd\Omega,
    \end{equation}
    因此$\phi[n+n_0]$是$\Phi(\Omega)\e^{-\im\Omega n_0}$的傅里叶系数。
    \begin{equation}
        \langle \tau_{n_0}x,\phi\rangle=\langle x,\tau_{-n_0}\phi\rangle=\langle X,\e^{-\im\Omega n_0}\Phi\rangle=\langle\e^{-\im \Omega n_0}X,\Phi\rangle,
    \end{equation}
    即$\DTFT(x[n-n_0])=\e^{-\im \Omega n_0}X$。
\end{proof}

\begin{proposition}[频移性质]
    \begin{equation}
        \DTFT(\e^{\im\Omega_0 n}x[n])=\tau_{\Omega_0}X.
    \end{equation}
\end{proposition}
\begin{proof}
    \begin{equation}
        \e^{\im\Omega_0 n}\phi[n]=\frac{1}{2\pi}\int_{-\pi}^{\pi}\Phi(\Omega)\e^{-\im(\Omega-\Omega_0) n}\,\dd\Omega=\frac{1}{2\pi}\int_{-\pi}^{\pi}\Phi(\Omega+\Omega_0)\e^{-\im\Omega n}\,\dd\Omega,
    \end{equation}
    因此$\e^{\im\Omega_0 n}\phi[n]$是$\Phi(\Omega+\Omega_0)$的傅里叶系数。
    \begin{equation}
        \langle \e^{\im\Omega_0 n}x,\phi\rangle=\langle x,\e^{\im\Omega_0 n}\phi\rangle=\langle X,\tau_{-\Omega_0}\Phi\rangle=\langle \tau_{\Omega_0}X,\Phi\rangle,
    \end{equation}
    即$\DTFT(\e^{\im\Omega_0 n}x[n])=\tau_{\Omega_0}X$。
\end{proof}

\begin{proposition}[共轭性质]
    \begin{equation}
        \DTFT(\overline{x[n]})=\overline{X^-}.
    \end{equation}
\end{proposition}
\begin{proof}
    \begin{equation}
        \overline{\phi[n]}=\frac{1}{2\pi}\overline{\int_{-\pi}^{\pi}\Phi(\Omega)\e^{-\im\Omega n}\,\dd\Omega}=\frac{1}{2\pi}\int_{-\pi}^{\pi}\overline{\Phi(\Omega)}\e^{\im\Omega n}\,\dd\Omega=\frac{1}{2\pi}\int_{-\pi}^{\pi}\overline{\Phi(-\Omega)}\e^{-\im\Omega n}\,\dd\Omega,
    \end{equation}
    因此$\overline{\phi[n]}$是$\overline{\Phi(-\Omega)}$的傅里叶系数。
    \begin{equation}
        \langle \overline{x},\phi\rangle=\overline{\langle x,\overline{\phi}\rangle}=\overline{\langle X,\overline{\Phi^-}\rangle}=\langle \overline{X^-},\Phi\rangle,
    \end{equation}
    最后一个等号用到了$\overline{X}^-=\overline{X^-}$。
\end{proof}
\begin{corollary}
    如果$x[n]$为实数序列，则$X$具有共轭对称性：$X=\overline{X^-}$。
\end{corollary}

\begin{proposition}[反转性质]
    \begin{equation}
        \DTFT(x^-[n])=X^-.
    \end{equation}
\end{proposition}
\begin{proof}
    \begin{equation}
        \phi[-n]=\frac{1}{2\pi}\int_{-\pi}^{\pi}\Phi(\Omega)\e^{\im\Omega n}\,\dd\Omega=\frac{1}{2\pi}\int_{-\pi}^{\pi}\Phi(-\Omega)\e^{-\im\Omega n}\,\dd\Omega,
    \end{equation}
    因此$\phi[-n]$是$\Phi(-\Omega)$的傅里叶系数。
    \begin{equation}
        \langle x^-,\phi\rangle=\langle x,\phi^-\rangle=\langle X,\Phi^-\rangle=\langle X^-,\Phi\rangle.
    \end{equation}
\end{proof}

\begin{proposition}[频域微分性质]\label{prop:DTFT_derivative}
    \begin{equation}
        \DTFT(nx[n])=\im X'.
    \end{equation}
\end{proposition}
\begin{proof}
    $\im \Phi'$的傅里叶系数为
    \begin{align}
        \frac{1}{2\pi}\int_{-\pi}^{\pi}\im \Phi'(\Omega)\e^{-\im \Omega n}\,\dd\Omega&=\frac{\im}{2\pi}\paren{\evalat{\Phi(\Omega)\e^{-\im \Omega n}}{-\pi}{\pi}-\int_{-\pi}^{\pi}\Phi(\Omega)(\e^{-\im \Omega n})'\,\dd\Omega} \nonumber \\
        &=\frac{n}{2\pi}\int_{-\pi}^{\pi}\Phi(\Omega)\e^{-\im\Omega n}\,\dd\Omega \nonumber \\
        &=n\phi[n].
    \end{align}
    因此
    \begin{equation}
        \langle nx,\phi\rangle=\langle x,n\phi\rangle=\langle X,\im\Phi'\rangle=\langle \im X',\Phi\rangle.
    \end{equation}
\end{proof}

\begin{theorem}[时域卷积定理]
    设$x\in\mathcal{T}(\mathbb{Z}),\ y,\phi\in\mathcal{S}(\mathbb{Z})$，则
    \begin{equation}
        \DTFT(x*y)=XY.
    \end{equation}
\end{theorem}
\begin{proof}
    \begin{equation}
        \langle(x*y),\phi\rangle=\langle x,(y^-)*\phi\rangle.
    \end{equation}
    令$\psi[n]=(y^-)*\phi[n]$，对应的频域函数为
    \begin{align}
        \Psi(\Omega)&=\sum_{k=-\infty}^{\infty}\psi[k]\e^{\im\Omega k} \nonumber \\
        &=\sum_{k=-\infty}^{\infty}\paren{\sum_{m=-\infty}^{\infty}y[m]\phi[k+m]}\e^{\im\Omega k} \nonumber \\
        &=\sum_{m=-\infty}^{\infty}y[m]\e^{-\im\Omega m}\sum_{n=-\infty}^{\infty}\phi[n]\e^{\im \Omega n} \qquad (n=k+m) \nonumber \\
        &=Y(\e^{\im\Omega})\Phi(\Omega).
    \end{align}
    因此
    \begin{equation}
        \langle(x*y),\phi\rangle=\langle x,(y^-)*\phi\rangle=\langle X,Y\Phi\rangle=\langle XY,\Phi\rangle.
    \end{equation}
\end{proof}

\begin{theorem}[频域卷积定理]
    \begin{equation}
        \DTFT(x[n]y[n])=X*Y.
    \end{equation}
\end{theorem}
\begin{proof}
    \begin{equation}
        \langle X*Y,\Phi\rangle=\langle X(u),\langle Y(v),\Phi(u+v)\rangle\rangle.
    \end{equation}
    令$\Psi(u)=\langle Y(v),\Phi(u+v)\rangle$，则其傅里叶系数为
    \begin{align}
        \psi[n]&=\frac{1}{2\pi}\int_{-\pi}^{\pi}\Psi(u)\e^{-\im un}\,\dd u \nonumber \\
        &=\braket{Y(v),\frac{1}{2\pi}\int_{-\pi}^{\pi}\Phi(u+v)\e^{-\im un}\,\dd u} \nonumber \\
        &=\braket{Y(v),\frac{1}{2\pi}\int_{-\pi}^{\pi}\Phi(\Omega)\e^{-\im (\Omega-v)n}\,\dd\Omega} \qquad (\Omega=u+v) \nonumber \\
        &=\langle Y(v),\phi[n]\e^{\im vn}\rangle \nonumber \\
        &=\frac{1}{2\pi}\phi[n]\int_{-\pi}^{\pi}Y(v)\e^{\im vn}\,\dd v \nonumber \\
        &=\phi[n]y[n].
    \end{align}
    因此
    \begin{align}
        \langle X*Y,\Phi\rangle&=\langle X(u),\langle Y(v),\Phi(u+v)\rangle\rangle \nonumber \\
        &=\langle X,\Psi\rangle \nonumber \\
        &=\langle x,\phi y\rangle \nonumber \\
        &=\langle xy,\phi\rangle.
    \end{align}
\end{proof}

由于$x[n],y[n]\in\mathcal{T}(\mathbb{Z})$时，$x[n]y[n]\in\mathcal{T}(\mathbb{Z})$，$X*Y$总是存在的，且\textbf{满足交换律和结合律}。

\begin{example}[正弦、余弦序列的离散时间傅里叶变换]\label{exmp:cos_DTFT}
    根据欧拉公式及DTFT的线性性可得：
    \begin{equation}
        \DTFT(\cos(\Omega_0 n))=\pi(\delta_{\Omega_0}+\delta_{-\Omega_0}),\quad \DTFT(\sin(\Omega_0 n))=-\pi\im(\delta_{\Omega_0}-\delta_{-\Omega_0}).
    \end{equation}
    当$\Omega_0/\pi$是无理数时，$\cos(\Omega_0 n)$和$\sin(\Omega_0 n)$不具有周期性，但是总存在有理数与$\Omega_0/\pi$充分靠近，并且频谱为线谱，因此可以说这种序列是\textbf{几乎周期}(almost periodic)的。我们在\refsub{power_signal}中讨论过类似的概念。
\end{example}

在上例中，令$\Omega_0=\pi$，则
\begin{equation}
    \DTFT(\cos(\pi n))=\DTFT((-1)^n)=\pi(\delta_\pi+\delta_{-\pi})=2\pi\delta_\pi.
\end{equation}
在这个意义下，它是振荡最快的序列。

\begin{example}[单位阶跃序列的离散时间傅里叶变换]
    \begin{align}
        & u[n]-u[n-1]=\delta[n] \nonumber \\
        \implies & \langle (1-\tau_1)u,\phi\rangle=\langle\delta,\phi\rangle \nonumber \\
        \implies & \langle(1-\e^{\im\Omega})U,\Phi\rangle=\langle 1,\Phi\rangle.
    \end{align}
    类似\refsub{normal_distribution}，作为一个分布方程，$(1-\e^{\im\Omega})U=1$并不能表明$U=1/(1-\e^{\im\Omega})$，由于
    \begin{equation}
        \frac{1}{1-\e^{\im\Omega}}\sim\frac{\im}{\Omega},\quad \Omega\to 0,
    \end{equation}
    $1/(1-\e^{\im\Omega})$在$[-\pi,\pi]$上不可积，也就无法与$1\in\mathcal{C}(\mathbb{T})$配对。我们需要更加仔细地考察$U$的性质。

    $(1-\e^{\im\Omega})U=1$确定了$U$在$\Omega\neq 0$处的值，为了使它能与$1$配对，我们只能认为$\langle U,\Phi\rangle$的意义是在主值意义下的积分，这样，$U$仍满足它的分布方程：
    \begin{equation}
        (1-\e^{\im\Omega})U=(1-\e^{\im\Omega})\pv\frac{1}{1-\e^{\im\Omega}}=1.
    \end{equation}
    由于$\Phi\in\mathcal{C}(\mathbb{T})$光滑，$1/(1-\e^{\im\Omega})$在$\Omega=0$附近的表现类似于奇函数$\im/\Omega$，该主值意义下的积分总是收敛的。

    主值积分不能描述$\Omega=0$处的$U$的行为，由于它是$1/(1-\e^{\im\Omega})$的极点，直观上可以使用$\delta$分布来处理这个问题。我们引入另一个分布方程：
    \begin{equation}
        u[n]+u[-n]=\begin{dcases}
            2, & \text{if } n=0,    \\
            1, & \text{if } n\neq 0
        \end{dcases}=1+\delta[n]\implies U+U^-=2\pi\delta+1.
    \end{equation}
    代入
    \begin{equation}
        U=\pv\frac{1}{1-\e^{\im\Omega}}+a\delta,
    \end{equation}
    得到$a=\pi$，因此：
    \begin{equation}
        U=\pv\frac{1}{1-\e^{\im\Omega}}+\pi\delta.
    \end{equation}
    对于z变换的收敛域边界恰为单位圆的序列，这样的做法是普适的。
\end{example}
\begin{example}
    本例探究$\delta'$的性质。
    \begin{equation}
        \langle \delta',\Phi\rangle=-\langle\delta,\Phi'\rangle=-\Phi'(0).
    \end{equation}
    与\cref{exmp:delta'}类似，$\delta'$的作用是取$0$处的导数的相反数。根据\cref{prop:DTFT_derivative}，有：
    \begin{equation}
        \DTFT(n)=\im\delta',
    \end{equation}
    这表明$\delta'$是$x[n]=-\im n$的DTFT。
\end{example}

利用分布的卷积的定义，容易得到
\begin{equation}
    \delta_a*\delta_b=\delta_{a+b},\quad \delta_a*X=\tau_a X.
\end{equation}
基于此，利用频域卷积定理可得：（注意周期函数的卷积带有系数$1/2\pi$）

\begin{example}[单边正弦、余弦序列的离散时间傅里叶变换]
    \begin{align}
        \DTFT(\cos(\Omega_0 n)u[n])&=\pi(\delta_{\Omega_0}+\delta_{-\Omega_0})*U \nonumber \\
        &=\pv\frac{1}{2}\paren{\frac{1}{1-\e^{\im(\Omega-\Omega_0)}}+\frac{1}{1-\e^{\im(\Omega+\Omega_0)}}}+\frac{\pi}{2}(\delta_{\Omega_0}+\delta_{-\Omega_0}) \nonumber \\
        &=\pv\frac{1-\e^{\im\Omega}\cos\Omega_0}{1-2\e^{\im\Omega}\cos\Omega_0+\e^{2\im\Omega}}+\frac{\pi}{2}(\delta_{\Omega_0}+\delta_{-\Omega_0}), \\
        \DTFT(\sin(\Omega_0 n)u[n])&=-\pi\im(\delta_{\Omega_0}-\delta_{-\Omega_0})*U \nonumber \\
        &=-\frac{\im}{2}\pv\paren{\frac{1}{1-\e^{\im(\Omega-\Omega_0)}}-\frac{1}{1-\e^{\im(\Omega+\Omega_0)}}}-\frac{\pi\im}{2}(\delta_{\Omega_0}-\delta_{-\Omega_0}) \nonumber \\
        &=-\pv\frac{\e^{\im\Omega}\sin\Omega_0}{1-2\e^{\im\Omega}\cos\Omega_0+\e^{2\im\Omega}}-\frac{\pi\im}{2}(\delta_{\Omega_0}-\delta_{-\Omega_0}).
    \end{align}
    与\cref{exmp:causual_cos_z}和\cref{exmp:casual_sin_z}对比，可以发现这两个序列的DTFT仍是z变换在单位圆上的取值，尽管单位圆并不在它们的z变换的收敛域内。
\end{example}

\subsection{分布的离散时间傅里叶变换的连续性}\label{sub:DTFT_continuous}

类似于\cref{thm:distribution_fourier_continuous}，我们有：
\begin{theorem}\label{thm:DTFT_continuous_1}
    设$\{X_k\}_{k=1}^\infty\subset\mathcal{D}(\mathbb{T}),\ X\in\mathcal{D}(\mathbb{T})$，如果$X_k\to X$，则$x_k[n]\to x[n],\ \forall n\in\mathbb{Z}$。
\end{theorem}
\begin{proof}
    \begin{equation}
        \lim_{k\to\infty}\langle x_k,\phi\rangle=\lim_{k\to\infty}\langle X_k,\Phi\rangle=\langle X,\Phi\rangle=\langle x,\phi\rangle,
    \end{equation}
    于是
    \begin{equation}
        \delta[n-m]\in\mathcal{S}(\mathbb{Z}),\ \forall m\in\mathbb{Z}\implies x_k[n]\to x[n],\quad \forall n\in\mathbb{Z}.
    \end{equation}
\end{proof}

然而，其逆命题需要额外的条件：
\begin{theorem}\label{thm:DTFT_continuous_2}
    设$\{x_k[n]\}_{k=1}^\infty\subset\mathcal{T}(\mathbb{Z}),\ x[n]\in\mathcal{T}(\mathbb{Z})$，如果$x_k[n]\to x[n],\ \forall n\in\mathbb{Z}$，并且$\{x_k[n]\}$具有\textbf{一致缓增性}，即：
    \begin{equation}
        \exists C>0,\ M\in\mathbb{N}^*,\quad |x_k[n]|\leq C(1+|n|)^M,\quad \forall k\in\mathbb{N}^*,\ n\in\mathbb{Z},
    \end{equation}
    则$X_k\to X$。
\end{theorem}
\begin{proof}
    设$\phi[n]\in\mathcal{S}(\mathbb{Z})$，$\Phi(\Omega)\in\mathcal{C}(\mathbb{T})$以$\phi[n]$为傅里叶系数，我们说明下式中的极限可以与求和交换：
    \begin{equation}
        \lim_{k\to\infty}\langle x_k,\phi\rangle=\lim_{k\to\infty}\sum_{n=-\infty}^{\infty}x_k[n]\phi[n].
    \end{equation}
    构造函数列
    \begin{equation}
        \{f_k:\mathbb{R}\to\mathbb{C}\}_{k=1}^\infty,\quad f_k(t)=\sum_{n=-\infty}^{\infty}x_k[n]\phi[n]\chi_{[n,n+1)}(t),
    \end{equation}
    显然$\{f_k\}$具有逐点极限
    \begin{equation}
        f:\mathbb{R}\to\mathbb{C},\quad f(t)=\sum_{n=-\infty}^{\infty}x[n]\phi[n]\chi_{[n,n+1)}(t).
    \end{equation}
    根据速降序列的定义，
    \begin{equation}
        \lim_{n\to\infty}Cn^2(1+|n|)^M\phi[n]=0,
    \end{equation}
    因此$\{Cn^2(1+|n|)^M\phi[n]\}_{k=1}^\infty$有上界$M'$，于是
    \begin{align}
        \abs{Cn^2(1+|n|)^M\phi[n]}&\leq \frac{M}{n^2}, \\
        \sum_{n=-\infty}^{\infty}\abs{Cn^2(1+|n|)^M\phi[n]}&\leq\sum_{n=-\infty}^{\infty}\frac{M}{n^2}<\infty.
    \end{align}
    由此构造出可积的控制函数：
    \begin{equation}
        g(t)=\sum_{n=-\infty}^{\infty}\abs{Cn^2(1+|n|)^M\phi[n]}\chi_{[n,n+1)}(t),\quad |f_k(t)|\leq g(t),\quad \forall k,\ \forall t.
    \end{equation}
    由控制收敛定理（\cref{thm:DCT}），可得：
    \begin{equation}
        \int_{\mathbb{R}}f\,\dd m=\lim_{k\to\infty}\int_{\mathbb{R}}f_k\,\dd m,
    \end{equation}
    即
    \begin{equation}
        \langle x,\phi\rangle=\sum_{n=-\infty}^{\infty}x[n]\phi[n]=\lim_{k\to\infty}\sum_{n=-\infty}^{\infty}x_k[n]\phi[n]=\lim_{k\to\infty}\langle x_k,\phi\rangle,
    \end{equation}
    因此
    \begin{equation}
        \langle X,\Phi\rangle=\langle x,\phi\rangle=\lim_{k\to\infty}\langle x_k,\phi\rangle=\lim_{k\to\infty}\langle X_k,\Phi\rangle.
    \end{equation}
\end{proof}

一致缓增性的条件不可以去除：
\begin{example}
    构造缓增序列
    \begin{equation}
        x_k[n]=\begin{dcases}
            k!, & n=k,    \\
            0,  & n\neq k.
        \end{dcases}
    \end{equation}
    固定$n$，当$k\to\infty$时，$x_k[n]\to 0$，因此$\{x_k[n]\}$逐点收敛到$0$序列$x[n]\equiv 0$。

    但是，取速降序列$\phi[n]=\displaystyle\frac{1}{n!}$，则
    \begin{equation}
        \langle x_k,\phi\rangle=\sum_{n=-\infty}^{\infty}x_k[n]\phi[n]=1,\quad \forall k\in\mathbb{N}^*,
    \end{equation}
    而$\langle X,\Phi\rangle=0$，这表明
    \begin{equation}
        \lim_{k\to\infty}\langle X_k,\Phi\rangle\neq \langle X,\Phi\rangle.
    \end{equation}
\end{example}

\subsection{序列的相关和谱密度}\label{sub:discrete_corr}

本小节仿照\refsec{Power_Signal}，先讨论能量信号的相关与能量谱，再对有限个复指数序列的叠加建立互功率、相关函数和功率谱。离散情形中，数字角频率以$2\pi$为周期，因此频率相差$2\pi$的整数倍时，必须视为同一频率。

\subsubsection{能量信号的相关与能量谱}

\begin{definition}[离散能量信号的相关]\label{def:discrete_correlation_energy}
    设$x,y\in l^2$，它们的\textbf{互相关}定义为
    \begin{equation}
        r_{xy}[l]=\sum_{n=-\infty}^{\infty}x[n+l]\overline{y[n]},\quad l\in\mathbb{Z}.
    \end{equation}
    当$x=y$时，称为$x$的\textbf{自相关}，记为$r_x[l]$。
\end{definition}

上述级数绝对收敛，因为
\begin{equation}
    \sum_{n=-\infty}^{\infty}|x[n+l]\overline{y[n]}|\leq\paren{\sum_{n=-\infty}^{\infty}|x[n+l]|^2}^{1/2}\paren{\sum_{n=-\infty}^{\infty}|y[n]|^2}^{1/2}=\|x\|_2\|y\|_2.
\end{equation}
特别地，定义两个信号的互能量为
\begin{equation}
    E_{xy}=\sum_{n=-\infty}^{\infty}x[n]\overline{y[n]},
\end{equation}
便有$r_{xy}[0]=E_{xy}$和$r_x[0]=E_x$。

\begin{proposition}[离散能量信号的相关的性质]
    设$x,y\in l^2$，则
    \begin{enumerate}
        \item $r_{xy}=x*\overline{y^-}$，并且$r_{xy}[l]=\overline{r_{yx}[-l]}$。
        \item 若$\tau_k x[n]=x[n-k]$，则$r_{\tau_k x,y}[l]=r_{xy}[l-k]=r_{x,\tau_{-k}y}[l]$。
        \item $|r_{xy}[l]|\leq\|x\|_2\|y\|_2$，特别地，$|r_x[l]|\leq r_x[0]=E_x$。
    \end{enumerate}
    若$x$为实序列，则$r_x$也是实序列，且为偶序列。
\end{proposition}
这些性质可由求和换元和上面的柯西—施瓦茨估计直接得到。

\begin{theorem}[离散能量信号的维纳-辛钦定理]\label{thm:discrete_energy_wiener_khinchin}
    设$x,y\in l^1\cap l^2$，其DTFT分别为$X$和$Y$，则
    \begin{equation}
        \DTFT(r_{xy})(\e^{\im\Omega})=X(\e^{\im\Omega})\overline{Y(\e^{\im\Omega})}.
    \end{equation}
\end{theorem}
\begin{proof}
    由于
    \begin{equation}
        \sum_{l=-\infty}^{\infty}\sum_{n=-\infty}^{\infty}|x[n+l]\overline{y[n]}|=\|x\|_1\|y\|_1<\infty,
    \end{equation}
    可以交换求和次序。令$m=n+l$，得到
    \begin{align}
        \DTFT(r_{xy})(\e^{\im\Omega})
        &=\sum_l\sum_n x[n+l]\overline{y[n]}\e^{-\im\Omega l} \nonumber \\
        &=\paren{\sum_m x[m]\e^{-\im\Omega m}}\paren{\sum_n\overline{y[n]}\e^{\im\Omega n}}
        =X\overline{Y}.
    \end{align}
\end{proof}

由DTFT的帕塞瓦尔定理（\cref{thm:DTFT_plancherel}），
\begin{equation}
    E_x=\frac{1}{2\pi}\int_{-\pi}^{\pi}|X(\e^{\im\Omega})|^2\,\dd\Omega,
    \quad
    E_{xy}=\frac{1}{2\pi}\int_{-\pi}^{\pi}X(\e^{\im\Omega})\overline{Y(\e^{\im\Omega})}\,\dd\Omega.
\end{equation}
这说明能量与互能量可以按数字角频率分配。

\begin{definition}[序列的能量谱密度]
    设$x,y\in l^1\cap l^2$，定义\textbf{能量谱密度}与\textbf{互能量谱密度}为
    \begin{equation}
        S_x^{\mathrm{en}}(\Omega)=\frac{1}{2\pi}|X(\e^{\im\Omega})|^2,
        \quad
        S_{xy}^{\mathrm{en}}(\Omega)=\frac{1}{2\pi}X(\e^{\im\Omega})\overline{Y(\e^{\im\Omega})}.
    \end{equation}
\end{definition}
这里保留系数$1/(2\pi)$，使谱密度在一个频率周期上的积分直接等于能量或互能量；上标$\mathrm{en}$用于与后面的功率谱区分。

\begin{example}[指数序列的相关与能量谱]
    设$x[n]=a^n u[n]$，其中$0<a<1$。当$l\geq 0$时，
    \begin{equation}
        r_x[l]=\sum_{n=0}^{\infty}a^{n+l}a^n=\frac{a^l}{1-a^2}.
    \end{equation}
    由实序列自相关的偶对称性，
    \begin{equation}
        r_x[l]=\frac{a^{|l|}}{1-a^2}.
    \end{equation}
    又因为$X(\e^{\im\Omega})=1/(1-a\e^{-\im\Omega})$，所以
    \begin{equation}
        S_x^{\mathrm{en}}(\Omega)=\frac{1}{2\pi(1-2a\cos\Omega+a^2)}.
    \end{equation}
    它在一个频率周期上的积分为$E_x=1/(1-a^2)$，其图像如下。
    \begin{figure}[H]
        \centering
        \includegraphics[width=0.6\textwidth]{exp_ESD}
        \caption{指数序列的能量谱密度}
    \end{figure}
\end{example}

\subsubsection{有限频率序列空间与平均内积}

对于序列$x[n]$，若时间平均极限存在，则记
\begin{equation}
    P_x=\lim_{M\to\infty}\frac{1}{2M+1}\sum_{n=-M}^{M}|x[n]|^2.
\end{equation}
满足$0<P_x<\infty$的信号称为\textbf{离散功率信号}。有限能量信号的平均功率为零，因为
\begin{equation}
    0\leq\frac{1}{2M+1}\sum_{n=-M}^{M}|x[n]|^2\leq\frac{E_x}{2M+1}\longrightarrow 0.
\end{equation}
因此，平均功率正而有限的信号必有无穷能量。不过，将零平均功率信号也纳入考虑，并不能保证全部信号构成线性空间。和\refsub{power_signal}一样，下面先选择一个具体的线性空间。

\begin{definition}[有限频率序列空间]
    将有限个复指数序列的线性组合构成的空间记为
    \begin{equation}
        \mathcal{P}_{\mathrm d}=\braces{x[n]=\sum_{k=1}^{K}a_k\e^{\im\Omega_k n}: K\in\mathbb{N}^*,\ a_k\in\mathbb{C},\ \Omega_k\in(-\pi,\pi]}.
    \end{equation}
\end{definition}

零序列也包含在$\mathcal{P}_{\mathrm d}$中；两个这样的序列相加或数乘，仍具有上述形式。每个$x\in\mathcal{P}_{\mathrm d}$都满足$|x[n]|\leq\sum_k|a_k|$，因而有界。

频率代表元取在$(-\pi,\pi]$内，是因为
\begin{equation}
    \e^{\im(\Omega+2\pi k)n}=\e^{\im\Omega n},\quad k,n\in\mathbb{Z}.
\end{equation}
因此，讨论频率表示时，约定先将频率按模$2\pi$归并，再合并同频率项。区间的两个端点$-\pi$和$\pi$代表同一个频率，不能重复计算。

由离散傅里叶反演公式（\cref{thm:DFT_inv}），任何周期为$N$的序列都能精确地表示为$N$个复指数序列之和，因而属于$\mathcal{P}_{\mathrm d}$。它的平均功率为
\begin{equation}
    P_x=\frac{1}{N}\sum_{n=0}^{N-1}|x[n]|^2.
\end{equation}
此外，$\mathcal{P}_{\mathrm d}$也包含非周期序列；例如，当$\Omega_0/(2\pi)$为无理数时，$\e^{\im\Omega_0 n}$不是周期序列。

\begin{proposition}[离散复指数的平均正交性]
    对任意$\Omega,\Omega'\in\mathbb{R}$，
    \begin{equation}
        \lim_{M\to\infty}\frac{1}{2M+1}\sum_{n=-M}^{M}\e^{\im(\Omega-\Omega')n}=
        \begin{dcases}
            1, & \Omega-\Omega'\in 2\pi\mathbb{Z},    \\
            0, & \Omega-\Omega'\notin 2\pi\mathbb{Z}.
        \end{dcases}
    \end{equation}
\end{proposition}
\begin{proof}
    若$\Omega-\Omega'\in 2\pi\mathbb{Z}$，每个求和项都等于$1$。否则，令$d=\Omega-\Omega'$，由有限等比数列求和公式，
    \begin{equation}
        \frac{1}{2M+1}\sum_{n=-M}^{M}\e^{\im dn}=\frac{\sin((2M+1)d/2)}{(2M+1)\sin(d/2)}\longrightarrow 0.
    \end{equation}
\end{proof}

设$x,y\in\mathcal{P}_{\mathrm d}$，在共同的有限频率集$\Lambda\subset(-\pi,\pi]$上写成
\begin{equation}
    x[n]=\sum_{\Omega\in\Lambda}a_\Omega\e^{\im\Omega n},\quad y[n]=\sum_{\Omega\in\Lambda}b_\Omega\e^{\im\Omega n}.
\end{equation}
未出现的频率对应的系数取零。由于求和有限，由平均正交性可得
\begin{equation}
    \lim_{M\to\infty}\frac{1}{2M+1}\sum_{n=-M}^{M}x[n]\overline{y[n]}=\sum_{\Omega\in\Lambda}a_\Omega\overline{b_\Omega}.
\end{equation}
将这个极限定义为\textbf{互功率}$P_{xy}$，并令
\begin{equation}
    (x,y)_{\pwr}=P_{xy}=\sum_{\Omega\in\Lambda}a_\Omega\overline{b_\Omega}.
\end{equation}
这是$\mathcal{P}_{\mathrm d}$上的内积：它对第一变元线性、对第二变元共轭线性，并且
\begin{equation}
    P_x=(x,x)_{\pwr}=\sum_{\Omega\in\Lambda}|a_\Omega|^2,\quad P_x=0\iff x=0.
\end{equation}
因此，$\mathcal{P}_{\mathrm d}$的非零元素都是功率信号，且
\begin{equation}
    \norm{x}_{\pwr}=\sqrt{P_x},\quad |P_{xy}|\leq\sqrt{P_xP_y}.
\end{equation}
这里先通过有限频率计算证明极限存在，再验证内积性质，不需要将全部功率信号视为一个空间。

\subsubsection{功率信号的相关与功率谱}

功率信号的相关采用时间平均。在$\mathcal{P}_{\mathrm d}$中，相关序列的存在性可直接由刚才的平均内积公式保证。

\begin{definition}[离散功率信号的相关]\label{def:discrete_correlation_power}
    设$x,y\in\mathcal{P}_{\mathrm d}$，它们的互相关定义为
    \begin{equation}
        r_{xy}[l]=\lim_{M\to\infty}\frac{1}{2M+1}\sum_{n=-M}^{M}x[n+l]\overline{y[n]}.
    \end{equation}
    当$x=y$时，称为自相关，记为$r_x[l]$。
\end{definition}

对每个固定的$l$，平移后的序列仍属于$\mathcal{P}_{\mathrm d}$，其频率$\Omega$处的系数为$a_\Omega\e^{\im\Omega l}$。由平均内积公式，
\begin{equation}
    r_{xy}[l]=\sum_{\Omega\in\Lambda}a_\Omega\overline{b_\Omega}\e^{\im\Omega l},\quad
    r_x[l]=\sum_{\Omega\in\Lambda}|a_\Omega|^2\e^{\im\Omega l}.
\end{equation}
这既证明了相关函数的存在性，也说明相关序列仍属于$\mathcal{P}_{\mathrm d}$。由此可得
\begin{equation}
    r_{xy}[0]=P_{xy},\quad r_x[0]=P_x,\quad
    r_{xy}[l]=\overline{r_{yx}[-l]},\quad
    |r_{xy}[l]|\leq\sqrt{P_xP_y}.
\end{equation}
特别地，$|r_x[l]|\leq P_x$；实序列的自相关为实偶序列。

为了将功率与频谱联系起来，将$x[n]$截短为长为$2M+1$的序列：
\begin{equation}
    x_M[n]=
    \begin{dcases}
        x[n], & -M\leq n\leq M, \\
        0,    & \text{其他}.
    \end{dcases}
\end{equation}
类似地，$y_M[n]$在$-M\leq n\leq M$时等于$y[n]$，其余位置为零。记$I_M=\{-M,\ldots,M\}$，以及$X_M=\DTFT(x_M)$、$Y_M=\DTFT(y_M)$。
截短序列只有有限个非零项，帕塞瓦尔定理给出
\begin{equation}
    \frac{1}{2M+1}\sum_{n=-M}^{M}|x[n]|^2=\int_{-\pi}^{\pi}\frac{|X_M(\e^{\im\Omega})|^2}{2\pi(2M+1)}\,\dd\Omega.
\end{equation}
左端趋于$P_x$，但右端被积函数有分布极限需要单独证明。仿照\refsub{power_spectrum}，先利用已知存在的相关序列定义功率谱，再证明截短谱的极限表示。

\begin{definition}[序列的功率谱密度]
    设$x,y\in\mathcal{P}_{\mathrm d}$，定义\textbf{功率谱密度}与\textbf{互功率谱密度}为
    \begin{equation}
        S_x=\frac{1}{2\pi}\DTFT(r_x),\quad
        S_{xy}=\frac{1}{2\pi}\DTFT(r_{xy}).
    \end{equation}
    这里的DTFT在周期分布的意义下理解。
\end{definition}

\begin{theorem}[有限频率序列的功率谱]\label{thm:discrete_PSD}
    设$x,y\in\mathcal{P}_{\mathrm d}$，其在共同频率集$\Lambda\subset(-\pi,\pi]$上的系数分别为$a_\Omega,b_\Omega$，则
    \begin{equation}
        S_x=\sum_{\Omega\in\Lambda}|a_\Omega|^2\delta_\Omega,\quad
        S_{xy}=\sum_{\Omega\in\Lambda}a_\Omega\overline{b_\Omega}\delta_\Omega.
    \end{equation}
    式中的$\delta_\Omega$沿用\cref{exmp:exp_DTFT}的周期分布记法。
\end{theorem}
\begin{proof}
    对相关序列的有限频率表达式逐项做DTFT，利用$\DTFT(\e^{\im\Omega l})=2\pi\delta_\Omega$，得到
    \begin{equation}
        \DTFT(r_{xy})=2\pi\sum_{\Omega\in\Lambda}a_\Omega\overline{b_\Omega}\delta_\Omega.
    \end{equation}
    除以$2\pi$即得互功率谱，令$y=x$即得自功率谱。
\end{proof}

每个频率周期内，各条自功率谱线的强度之和为$P_x$，互功率谱线的强度之和为$P_{xy}$。

注意，原序列的DTFT是
\begin{equation}
    X=2\pi\sum_{\Omega\in\Lambda}a_\Omega\delta_\Omega.
\end{equation}
因此，若直接将DTFT写成$X=\sum_{\Omega\in\Lambda}c_\Omega\delta_\Omega$，则$c_\Omega=2\pi a_\Omega$，相应功率谱为
\begin{equation}
    S_x=\sum_{\Omega\in\Lambda}\frac{|c_\Omega|^2}{4\pi^2}\delta_\Omega.
\end{equation}

对于周期为$N$的序列，若一周期的DFT为$X_N[k]$，则离散傅里叶展开系数为$X_N[k]/N$。令$\Omega_k$是在$(-\pi,\pi]$内与$2\pi k/N$同余的频率，即得
\begin{equation}
    S_x=\sum_{k=0}^{N-1}\frac{|X_N[k]|^2}{N^2}\delta_{\Omega_k},\quad
    P_x=\sum_{k=0}^{N-1}\frac{|X_N[k]|^2}{N^2}=\frac{1}{N}\sum_{n=0}^{N-1}|x[n]|^2.
\end{equation}

\subsubsection{截短功率谱的分布极限}

\begin{theorem}[离散功率信号的维纳-辛钦定理]\label{thm:deiscrete_winner_khinchin}
    设$x,y\in\mathcal{P}_{\mathrm d}$，$x_M,y_M$为上述截短序列，其DTFT分别为$X_M,Y_M$。则在周期分布的意义下，
    \begin{equation}
        \lim_{M\to\infty}\frac{X_M\overline{Y_M}}{2M+1}=\DTFT(r_{xy})=2\pi S_{xy}.
    \end{equation}
    特别地，
    \begin{equation}
        \lim_{M\to\infty}\frac{|X_M|^2}{2\pi(2M+1)}=S_x.
    \end{equation}
\end{theorem}
\begin{proof}
    考察截短序列的互相关除以窗口长度所得的序列
    \begin{align}
        q_M[l]
        &=\frac{1}{2M+1}\sum_{n=-\infty}^{\infty}x_M[n+l]\overline{y_M[n]} \nonumber \\
        &=\frac{1}{2M+1}\sum_{n\in I_M\cap(I_M-l)}x[n+l]\overline{y[n]}.
    \end{align}
    必须同时保留两个截短窗口。固定$l$，当$M>|l|$时，与
    \begin{equation}
        c_M[l]=\frac{1}{2M+1}\sum_{n=-M}^{M}x[n+l]\overline{y[n]}
    \end{equation}
    相比，$q_M[l]$少了$|l|$项。故
    \begin{equation}
        |q_M[l]-c_M[l]|\leq\frac{|l|}{2M+1}\|x\|_\infty\|y\|_\infty\longrightarrow 0.
    \end{equation}
    已知$c_M[l]\to r_{xy}[l]$，因此$q_M[l]\to r_{xy}[l]$。同时，对所有$M,l$都有
    \begin{equation}
        |q_M[l]|\leq\|x\|_\infty\|y\|_\infty.
    \end{equation}
    这给出了\cref{thm:DTFT_continuous_2}所需的一致缓增性，因而$\DTFT(q_M)\to\DTFT(r_{xy})$。

    也可以直接核对分布收敛。任取$\Phi\in\mathcal{C}(\mathbb{T})$，其傅里叶系数$\phi[l]$速降，故$\sum_l|\phi[l]|<\infty$。利用与$M$无关的可求和控制项$\|x\|_\infty\|y\|_\infty|\phi[l]|$，可将极限与求和交换：
    \begin{equation}
        \lim_{M\to\infty}\sum_l q_M[l]\phi[l]=\sum_l r_{xy}[l]\phi[l].
    \end{equation}
    由DTFT的配对定义，这正是上述周期分布收敛。

    最后，对有限支集序列$x_M,y_M$应用\cref{thm:discrete_energy_wiener_khinchin}，得到
    \begin{equation}
        \DTFT(q_M)=\frac{X_M\overline{Y_M}}{2M+1}\longrightarrow\DTFT(r_{xy})=2\pi S_{xy}.
    \end{equation}
\end{proof}

\subsubsection{功率谱的计算例子}

\begin{example}[余弦序列的功率谱]
    设$x[n]=A\cos(\Omega_0n+\theta)$，其中$A\in\mathbb{R}$、$0<\Omega_0<\pi$。由
    \begin{equation}
        x[n]=\frac{A}{2}\e^{\im\theta}\e^{\im\Omega_0n}+\frac{A}{2}\e^{-\im\theta}\e^{-\im\Omega_0n},
    \end{equation}
    得到
    \begin{equation}
        P_x=\frac{A^2}{2},\quad
        r_x[l]=\frac{A^2}{2}\cos(\Omega_0l),\quad
        S_x=\frac{A^2}{4}(\delta_{\Omega_0}+\delta_{-\Omega_0}).
    \end{equation}
    这一结论不要求$\Omega_0/(2\pi)$为有理数，因此也适用于非周期余弦序列。
\end{example}

为了与截短相关的证明对应，取$x[n]=\cos(\Omega_0n)$、$0<\Omega_0<\pi$，记
\begin{equation}
    r_M[l]=\sum_{n=-\infty}^{\infty}x_M[n+l]x_M[n].
\end{equation}
当$|l|>2M$时，两个窗口不相交，$r_M[l]=0$；当$|l|\leq 2M$时，交集包含$2M+1-|l|$个整数。由积化和差及有限等比数列求和，
\begin{equation}
    r_M[l]=\frac{1}{2}\sbr{\frac{\sin((2M+1-|l|)\Omega_0)}{\sin\Omega_0}+(2M+1-|l|)\cos(\Omega_0l)}.
\end{equation}
于是对每个固定的$l$，
\begin{equation}
    \frac{r_M[l]}{2M+1}\longrightarrow\frac{1}{2}\cos(\Omega_0l)=r_x[l],
\end{equation}
与前面的结论一致。
